\documentclass[11pt,a4paper]{article}
\usepackage{curso}
\begin{document}
\portada{Lección}{2.5}{El teorema de fluctuación--disipación}

\noindent\textbf{Unidad 2.}\quad\textbf{Objetivos:} enunciar el TFD; verificarlo en el OU; reconocer Einstein y Nyquist como
casos particulares y la violación fuera del equilibrio. \textbf{Prerrequisitos:} 2.2--2.4, 1.6.

\section{Enunciado}
Sea un sistema en equilibrio a temperatura $T$ y una perturbación pequeña $-F\,x$ en la energía, encendida en $t=0$.
Con $C(t)=\langle\delta x(0)\delta x(t)\rangle_{\text{eq}}$ la autocorrelación sin perturbar:
\begin{resultado}[ (TFD clásico)]
\begin{equation}\langle x\rangle_F(t)-\langle x\rangle_{\text{eq}}=\beta F\,[C(0)-C(t)]+O(F^2),\qquad\chi(t)=-\beta\,\dot C(t)\ (t>0).\end{equation}
En frecuencia: $S_x(\omega)=\frac{2\kB T}{\omega}\mathrm{Im}\,\chi(\omega)$.
\end{resultado}
\textbf{Interpretación:} la respuesta a un empujón pequeño es la relajación de una fluctuación espontánea (hipótesis de
regresión de Onsager). La misma fricción que disipa el empujón genera las fluctuaciones \cite{peliti,seifert}.

\section{Verificación en el OU}
Para $dx=(-\kappa x+F)dt+\sqrt{2\kB T}\,dW$ ($\mu=1$): $C(t)=\frac{\kB T}{\kappa}e^{-\kappa t}$ y
$\langle x\rangle_F(t)=\frac{F}{\kappa}(1-e^{-\kappa t})=\beta F[C(0)-C(t)]$ exactamente. La animación lo comprueba
numéricamente \evS: $C(t)$ medida en una trayectoria de longitud 4000 y $\langle x\rangle_F$ promediada sobre $10^4$
partículas con $F=0{,}2$, $\kappa=1$, $\kB T=0{,}5$; ambas coinciden dentro del error estadístico ($\sim0{,}007$).

\section{Casos particulares}
\begin{itemize}[nosep]
\item \textbf{Einstein:} partícula libre, respuesta $=$ movilidad $\mu$, fluctuación $=$ difusión: $D=\mu\kB T$.
\item \textbf{Nyquist:} circuito, respuesta $=$ impedancia, fluctuación $=$ ruido de voltaje: $S_V=4\kB TR$.
\item \textbf{Equipartición:} respuesta estática $\chi_0=1/\kappa$ y $C(0)=\kB T/\kappa$: $C(0)=\kB T\chi_0$.
\end{itemize}

\section{Fuera del equilibrio}
En estados forzados (NESS) o con ruido no térmico, el TFD no vale. Se define una temperatura efectiva
$\kB T_{\text{ef}}(t)=\dot C(t)/(-\chi(t))$; si $T_{\text{ef}}\neq T$, hay forzamiento o ruido de otro origen. Herramienta de
diagnóstico para hardware termodinámico (ruido $1/f$, ruido de disparo, desajustes; Unidad 5.6).

\begin{ejemplo}
Pinza óptica: $\kappa=10^{-6}$ N/m, $T=300$ K: $\langle x^2\rangle=\kB T/\kappa\approx4{,}1\times10^{-15}$ m$^2$ ($x_{\text{rms}}\approx64$ nm).
Sin aplicar fuerza, el TFD predice la respuesta estática a $F=1$ pN: $\langle x\rangle=\beta F\,C(0)=F/\kappa=1\ \mu$m.
Recíprocamente, midiendo $x_{\text{rms}}$ se calibra $\kappa=\kB T/\langle x^2\rangle$.
\end{ejemplo}

\begin{aplicacion}
Calibración de trampas y sensores; termometría de ruido $T=\langle V^2\rangle/(4\kB R\Delta f)$; diagnóstico de chips; compromiso
de diseño: más disipación implica más ruido.
\end{aplicacion}

\begin{error}
\begin{itemize}[nosep]
\item Fuera del régimen lineal o del equilibrio no aplica.
\item $\chi(t)$ (impulso) $\neq$ respuesta a un escalón $\int_0^t\chi$.
\item El factor $\beta$; versión cuántica si $\hbar\omega\gtrsim\kB T$.
\end{itemize}
\end{error}

\section{Resumen}
Respuesta lineal $=$ relajación de fluctuaciones; $\langle x\rangle_F=\beta F[C(0)-C(t)]$; Einstein, Nyquist y equipartición
como casos; violación $\Rightarrow$ temperatura efectiva.

\begin{thebibliography}{9}\small
\bibitem[2]{peliti} L. Peliti y S. Pigolotti, \emph{Stochastic Thermodynamics: An Introduction}, 2021.
\bibitem[11]{seifert} U. Seifert, Stochastic thermodynamics, fluctuation theorems and molecular machines, \emph{Rep. Prog. Phys.} 75, 126001 (2012).
\bibitem[1]{gardiner} C. W. Gardiner, \emph{Handbook of Stochastic Methods}.
\end{thebibliography}
\end{document}
